\hypertarget{link_to_timer}{}{\Large \bf timer*}

{\it Timer with varying delay. The latter is a piecewise linear function of the monitored variable

If $x_i$ is smaller than a threshold $v_1$, the output $x_j$ is equal to zero. Otherwise, $x_j$ changes from zero to one at time $t^{\star}+\tau(x_i)$ where $t^{\star}$ is the time at which the input $x_i$ became larger than $v_1$ and the delay $\tau(x_i)$ varies with $x_i$ according to a piecewise linear characteristic involving $n$ points (see diagram below). 

Separate blocks exist for $n=1,2,3,4$ and $5$.}
\begin{center}
\includegraphics{codegen/timer.pdf}
\end{center}

\vspace*{-4ex}
Syntax : \hspace*{10ex}\parbox[t]{\textwidth}{\sf 
\& timer1\\
name of variable $x_i$\\
name of variable $x_j$\\
data name, parameter name or math expression for $v_1$\\
data name, parameter name or math expression for $T_1$
}

~

\hspace*{-8ex}\parbox[t]{100mm}{\sf 
\& timer2\\
name of variable $x_i$\\
name of variable $x_j$\\
data name, parameter name or math expression for $v_1$\\
data name, parameter name or math expression for $T_1$\\
data name, parameter name or math expression for $v_2$\\
data name, parameter name or math expression for $T_2$
}\parbox[t]{95mm}{\sf 
\& timer3\\
name of variable $x_i$\\
name of variable $x_j$\\
data name, parameter name or math expression for $v_1$\\
data name, parameter name or math expression for $T_1$\\
data name, parameter name or math expression for $v_2$\\
data name, parameter name or math expression for $T_2$\\
data name, parameter name or math expression for $v_3$\\
data name, parameter name or math expression for $T_3$
}

~

\hspace*{-8ex}\parbox[t]{100mm}{\sf 
\& timer4\\
name of variable $x_i$\\
name of variable $x_j$\\
data name, parameter name or math expression for $v_1$\\
data name, parameter name or math expression for $T_1$\\
data name, parameter name or math expression for $v_2$\\
data name, parameter name or math expression for $T_2$\\
data name, parameter name or math expression for $v_3$\\
data name, parameter name or math expression for $T_3$\\
data name, parameter name or math expression for $v_4$\\
data name, parameter name or math expression for $T_4$
}\parbox[t]{95mm}{\sf 
\& timer5\\
name of variable $x_i$\\
name of variable $x_j$\\
data name, parameter name or math expression for $v_1$\\
data name, parameter name or math expression for $T_1$\\
data name, parameter name or math expression for $v_2$\\
data name, parameter name or math expression for $T_2$\\
data name, parameter name or math expression for $v_3$\\
data name, parameter name or math expression for $T_3$\\
data name, parameter name or math expression for $v_4$\\
data name, parameter name or math expression for $T_4$\\
data name, parameter name or math expression for $v_5$\\
data name, parameter name or math expression for $T_5$
}

~ \hrule ~

Internal states : $x_1$

Discrete variables : $z \in \{-1,0,1\}$

Equations :
\begin{equation*}
0 =
\begin{cases}
x_j \;\;\;\;\;\;\;\;\;\; \mbox{if} \;\; z \in \{-1,0 \} \\
x_j-1 \;\;\;\;\;\;\;\;\;\; \mbox{if} \;\; z = 1
\end{cases} 
\end{equation*}
\begin{equation*}
\begin{cases}
0 & = x_1 \;\;\;\;\;\;\;\;\;\; \mbox{if} \;\; z = -1 \\
\dot{x_1} & =  1 \;\;\;\;\;\;\;\;\;\; \mbox{if} \;\; z = 0 \\
\dot{x_1} & =  0 \;\;\;\;\;\;\;\;\;\; \mbox{if} \;\; z = 1
\end{cases}
\end{equation*}

\pagebreak
Discrete transitions : 

\hspace*{20ex}\begin{minipage}[t]{6cm}
\begin{algorithmic} 
\IF {$z = -1$}
  \IF {$x_i \ge v_1$} 
    \STATE $z \gets 0$
  \ENDIF
\ELSE
  \IF {$x_i < v_1$} 
     \STATE $z \gets -1$
  \ENDIF
\ENDIF
\IF {$z = 0$}
  \IF {$x_1 \ge \tau(x_i)$}
    \STATE $z \gets 1$
  \ENDIF
\ENDIF
\end{algorithmic}
\end{minipage}


\vspace*{2ex}
Initialization of internal states : ~~~ $x_1 \gets 0$

\vspace*{2ex}
Initialization of discrete variables : ~~~ $z \gets -1$


The $v_i$ values must be increasing, but two consecutive values may be equal, i.e. $v_1 \le v_2 \le v_3 \le \ldots \le v_{n-1} \le v_n$.

The piecewise linear characteristic is typically used to approximate an inverse-time characteristic, in which case the $T$ values are decreasing, i.e. $T_1 \ge T_2 \ge T_3 \ge \ldots \ge T_{n-1} \ge T_n$. Nevertheless, non decreasing values are also allowed.

If the initial value of $x_i$ is larger than $v_1$, $x_j$ will change to one after the time $\tau(x_i)$, unless $x_i$ decreases below $v_1$ before the delay $\tau$ is elapsed.